\documentclass[10pt,a4paper]{amsart}
\usepackage[margin=19mm]{geometry}
\usepackage{amsmath,amssymb,mathtools}
\usepackage[scr=rsfs,cal=euler]{mathalfa}
\usepackage{xfrac}
\usepackage[unicode,hidelinks,bookmarks=false]{hyperref}
\newcommand{\ie}{i.e.\ }
\newcommand{\cf}{cf.\ }
\newcommand{\resp}{respectively\ }
\newcommand{\Subsection}{\subsection}
\providecommand{\nobreakdash}{\nobreak\hbox{-}}
\setlength{\emergencystretch}{2em}
\multlinegap0pt
\usepackage{bm}
\usepackage{bbm}
\usepackage{dsfont}
\usepackage{xspace}
\usepackage{cancel}
\usepackage{mathtools}
\usepackage{amsthm}
\makeatletter
\@ifundefined{theorem}{\newtheorem{theorem}{Theorem}}{}
\@ifundefined{lemma}{\newtheorem{lemma}[theorem]{Lemma}}{}
\@ifundefined{remark}{\newtheorem*{remark}{Remark}}{}
\makeatother
\newcommand\modd{\operatorname{mod}}
\title[A Generalisation of Sperner's Theorem Using Weighted Chains]{A Generalisation of Sperner's Theorem Using Weighted Chains}
\author{Ya\"el Dillies, Matthew Johnson, Aleksandra Kowalska}
\date{Source: arXiv:2509.26493v1 (CC BY 4.0)}
\begin{document}
\maketitle

\noindent\emph{Source and licence.} A Generalisation of Sperner's Theorem Using Weighted Chains. Version arXiv:2509.26493v1 (CC BY 4.0), distributed under the Creative Commons Attribution 4.0 International licence, as recorded in the arXiv licence metadata for this exact version and shown on the abstract page https://arxiv.org/abs/2509.26493v1. Full rights notice: licenses/arxiv-2509.26493-CC-BY-4.0.txt.\par\smallskip

\noindent\textit{Editorial scope.} Counted unit: the Remark whose source label is \texttt{None} in the article (source0.tex lines 1016--1018, source label \texttt{None}), delivered together with the article's own complete original proof of it (source0.tex lines 1020--1099, one \texttt{proof} environment of 80 lines). No mathematics is written, repaired or re-derived: statement and proof are the article's own text line for line. Required context retained uncounted: eq\_step\_1 (lines 607--610); eq\_step\_2 (lines 643--646); s\_closed\_form (lines 669--675); eq\_step\_1\_prime (lines 713--716); s\_closed\_form\_prime (lines 723--729). Every definition, display and label of those lines is retained unchanged. Omitted from this package and not redistributed: 1--606, 611--642, 647--668, 676--712, 717--722, 730--1015, 1019--1019, 1100--1204. Nothing inside a retained excerpt is omitted or altered: every retained body is delivered exactly as the article's lines are written, so no omission receipt of any kind accompanies this package. Citations: the retained excerpts carry no citation command at all, so no bibliography entry of the article is transcribed and no reference is redistributed. Reference adaptation: none was needed and none was made; every reference inside the delivered unit resolves inside it, so no unresolved reference or citation is produced. Printed slips kept literal: no printed word, symbol, hypothesis or number of the counted statement or of its proof was altered. Layout and class: the article's own class is \texttt{amsart} with packages including babel, xcolor, inputenc, amsmath, amsfonts, amsthm, amscd, amssymb, fullpage, upgreek, enumerate, comment, siunitx, enumitem; the wrapper is the lane's pinned \texttt{amsart} editorial wrapper at 10pt on A4, which restates the theorem-like environments the retained text needs and adds the article's own macro definitions that the retained text needs (\texttt{\textbackslash modd}), with extra wrapper packages (bm, bbm, dsfont, xspace, cancel, mathtools). The delivered numbering is the unit's own. Assets: no retained span contains a figure environment, a graphics-inclusion command, a PDF-page inclusion, a table or a TikZ picture; no image file is copied into this package, the package declares no asset dependency and no image file is redistributed with it. Licence: version arXiv:2509.26493v1 of this article is distributed under the Creative Commons Attribution 4.0 International licence, as recorded in the arXiv licence metadata for this exact version and shown on the abstract page https://arxiv.org/abs/2509.26493v1. A full rights notice is delivered at licenses/arxiv-2509.26493-CC-BY-4.0.txt. Characters: every retained excerpt is pure ASCII, so the delivered engine, XeLaTeX with its default Latin Modern text font, reports no missing character.
\begin{equation}
\label{eq_step_1}
    W(a, c) = S(a, c) - S(a+1, c-1) + R(a+1+k, c-1-k) + R(a+1+k, c-k).
\end{equation}

\begin{equation}
\label{eq_step_2}
    R(d, a, c)=S(d, a, c) - S(d, a+1, c-1) + R(d+1, a+1+k, c-1-k)+ R(d+1, a+1+k, c-k).
\end{equation}

\begin{lemma}
\label{s_closed_form}
    For any $d, a, c \geqslant 0$ we have:
    \begin{equation*}
        S(d, a, c) = \binom{n}{c} \cdot \binom{n-c-d-1}{n-a-c}
    \end{equation*}
\end{lemma}

\begin{equation}
\label{eq_step_1_prime}
    W(a, c) = S'(0, a, c) - S'(0, a-1, c+1) + R'(0, a-1-k, c+1+k) + R'(0, a-k, c+1+k).
\end{equation}

\begin{lemma}
\label{s_closed_form_prime}
    We have
    \begin{equation*}
        S'(d, a, c) = \binom{n}{a} \binom{n-a-d-1}{n-a-c}.
    \end{equation*}
\end{lemma}

\begin{remark}
    The proof of this lemma mostly involves binomial coefficient manipulations, which might seem a bit unmotivated. Plotting contributions from different types on a staircase diagram was very helpful for arriving at it.
\end{remark}

\begin{proof}
    We will show that for $(A, C)$ a lower type, we have
    \begin{multline}
    \label{lower_ac_part}
        U(A, C) = \sum_{\substack{c \equiv C \ (\modd k+1), \\ c \leqslant C}} \sum_{j \geqslant 0} \binom{l_C(c)}{j} \binom{n}{c+j} \binom{n-c-j-l_C(c)-1}{B-j}\\ - \sum_{\substack{c \equiv C-1 \ (\modd k+1), \\ c \leqslant C-1}} \sum_{j \geqslant 0} \binom{l_{C-1}(c)}{j} \binom{n}{c+j} \binom{n-c-j-l_{C-1}(c)-1}{B-j}
    \end{multline}
    and that for $(A, C)$ an upper type, we have
    \begin{multline}
    \label{upper_ac_part}
        U(A, C) = \sum_{\substack{c \equiv C \ (\modd k+1), \\ c \geqslant C}} \sum_{j \geqslant 0} \binom{l_{C}(c)-1}{j} \binom{n}{c+j} \binom{n-c-j-(l_{C}(c)-1)-1}{B-j}\\ - \sum_{\substack{c \equiv C+1 \ (\modd k+1), \\ c \geqslant C+1}} \sum_{j \geqslant 0} \binom{l_{C+1}(c)-1}{j} \binom{n}{c+j} \binom{n-c-j-(l_{C+1}(c)-1)-1}{B-j},
    \end{multline}
    which together (when we substitute $A-k, C+k$ for $A, C$ in the second one) imply the lemma.

    Let us first show Equation \ref{lower_ac_part}. To do so, we will mostly use Equations \ref{eq_step_1} and \ref{eq_step_2}. For easier manipulations, let us denote
    \begin{equation*}
        D(d, A, C) := S(d, A, C) - S(d, A+1, C-1).
    \end{equation*}
    For any $I$, we have:
    \begin{multline*}
        U(A, C) = \sum_{0 \leqslant i \leqslant I} \sum_{0 \leqslant j \leqslant i} \binom{i}{j} D(i, A+i(1+k), C-i(1+k)+j)\\
        + \sum_{0 \leqslant j \leqslant I} \binom{I}{j} \bigg( R\Big(I+1, A+(I+1)(1+k), C-(I+1)(1+k)+j\Big)\\ + R\Big(I+1, A+(I+1)(1+k), C-(I+1)(1+k)+j+1\Big) \bigg)
    \end{multline*}
    (for $I=0$, this is simply Equation \ref{eq_step_1}, and for greater $I$s it follows by induction and by applying \ref{eq_step_2}).

    We note that since for $j > i$ we have $\binom{i}{j}=0$, we can change the inner summation from range $0 \leqslant j \leqslant i$ to range $0 \leqslant j$.
    
    Further noticing that for $I$ sufficiently large, $C- Ik < 0$ and for any $j \leqslant I$ we have
    \begin{multline*}
        R\Big(I+1, A+(I+1)(1+k), C-(I+1)(1+k)+j\Big)\\ = R\Big(I+1, A+(I+1)(1+k), C-(I+1)(1+k)+j+1\Big) = 0,
    \end{multline*}
    we get that
    \begin{equation*}
        U(A, C) = \sum_{i \geqslant 0} \sum_{j \geqslant 0} \binom{i}{j} D(i, A+i(1+k), C-i(1+k)+j).
    \end{equation*}
    After applying Lemma \ref{s_closed_form}, we get that
    \begin{multline*}
        U(A, C) = \sum_{i \geqslant 0} \sum_{j \geqslant 0} \binom{i}{j} \binom{n}{C-i(1+k)+j} \binom{n-C+i(k+1)-j-i-1}{B-j}\\ - \sum_{i \geqslant 0} \sum_{j \geqslant 0} \binom{i}{j} \binom{n}{C-1-i(1+k)+j} \binom{n-C+i(k+1)-j-i}{B-j}.
    \end{multline*}
    Changing from summation over $i$ to summation over $c \equiv C \ (\modd k+1)$ (so $i$ becomes $\frac{C-c}{k+1}=l_C(c)$), we further get that
    \begin{multline*}
        U(A, C) = \sum_{\substack{c \equiv C \ (\modd k+1), \\ c \leqslant C}} \sum_{j \geqslant 0} \binom{l_C(c)}{j} \binom{n}{c+j} \binom{n-c-j-l_C(c)-1}{B-j}\\ - \sum_{\substack{c \equiv C-1 \ (\modd k+1), \\ c \leqslant C-1}} \sum_{j \geqslant 0} \binom{l_{C-1}(c)}{j} \binom{n}{c+j} \binom{n-c-j-l_{C-1}(c)-1}{B-j},
    \end{multline*}
    as desired.

    Now, we need to prove Equation \ref{upper_ac_part}.
    Using Equations \ref{eq_step_1_prime}, \ref{eq_step_2} and Lemma \ref{s_closed_form_prime}, analogously to our considerations above, we get that for $(A, C)$ an upper type we have
    \begin{multline*}
        U(A, C) = \sum_{i \geqslant 0} \sum_{j \geqslant 0} \binom{i}{j} \Big( S'(i, A-i(1+k)+j, C+i(1+k)) - S'(i, A-i(1+k)+j-1, C+i(1+k)+1) \Big)\\
        = \sum_{\substack{a \equiv A \ (\modd k+1), \\ a \leqslant A}} \sum_{j \geqslant 0}  \binom{l_A(a)}{j} \bigg( \binom{n}{a+j}\binom{n-a-j-i-1}{B-j} - \binom{n}{a+j-1}\binom{n-a-j-i}{B-j} \bigg)\\
        = \sum_{\substack{a \equiv A \ (\modd k+1), \\ a \leqslant A}} \sum_{j \geqslant 0}  \binom{l_A(a)}{j}\binom{n}{a+j}\binom{n-a-j-i-1}{B-j}\\ - \sum_{\substack{a \equiv A-1 \ (\modd k+1), \\ a \leqslant A}} \sum_{j \geqslant 0}  \binom{l_A(a)}{j}\binom{n}{a+j}\binom{n-a-j-i-1}{B-j}.
    \end{multline*}
    Hence, it suffices to show that
    \begin{multline}
    \label{what_we_need_lemma}
        \sum_{\substack{a \equiv A \ (\modd k+1), \\ a \leqslant A}} \sum_{i \geqslant 0}  \binom{l_A(a)}{i}\binom{n}{a+i}\binom{n-a-i-l_A(a)-1}{B-i}\\ = \sum_{\substack{c \equiv C \ (\modd k+1), \\ c \geqslant C}} \sum_{j \geqslant 0} \binom{l_{C}(c)-1}{j} \binom{n}{c+j} \binom{n-c-j-(l_{C}(c)-1)-1}{B-j}
    \end{multline}
    (where we changed the summation over $j$ on the left-hand side into a summation over $i$, as this will make our calculations clearer later). 
    First, we note that for $c := n-B-a$, we have
    \begin{multline*}
        \binom{l}{i} \binom{n}{a+i}\binom{n-a-i-l-1}{B-i}=\binom{l}{i} \binom{n}{a+i}\binom{n-a-i-l-1}{n-c-a-i}\\
        = \binom{l}{i} \binom{n}{a+i} \sum_{j \geqslant 0} \binom{-l-1}{j} \binom{n-a-i}{n-c-a-j}\\
        = \binom{n}{a+i} \sum_{j \geqslant 0} (-1)^j \binom{l}{i}\binom{l+j}{j} \binom{n-a-i}{n-c-a-j}
        = \sum_{j \geqslant 0} (-1)^j \binom{l+j}{i, j}\binom{n}{a+i, c+j}.
    \end{multline*}
    Hence, the left-hand side of Equation \ref{what_we_need_lemma} equals
    \begin{equation}
    \label{lhs_first_transformation}
        \sum_{\substack{a \equiv A \ (\modd k+1), \\ a \leqslant A}} \sum_{i \geqslant 0} \sum_{j \geqslant 0}  (-1)^j \binom{l_A(a)+j}{i + j}\binom{i+j}{i}\binom{n}{a+i, c+j}.
    \end{equation}
    Since $c=n-B-a$, we can change the summation from $\sum_{\substack{a \equiv A \ (\modd k+1), \\ a \leqslant A}}$ to $\sum_{\substack{c \equiv C \ (\modd k+1), \\ c \geqslant C}}$, where $l_A(a)=\frac{A-a}{k+1}=\frac{c-C}{k+1}=-l_C(c)$. Hence, the left-hand side of Equation \ref{what_we_need_lemma} equals
    \begin{equation*}
        \sum_{\substack{c \equiv C \ (\modd k+1), \\ c \geqslant C}} \sum_{i \geqslant 0} \sum_{j \geqslant 0} (-1)^j \binom{-l_C(c)+j}{i+j}\binom{i+j}{i}\binom{n}{a+i, c+j}.
    \end{equation*}
    Using the fact that for any $p, q$ positive we have $\binom{p-q}{p}(-1)^p=\binom{q-1}{p}$, we can transform it to
    \begin{equation}
    \label{expression_lhs}
        \sum_{\substack{c \equiv C \ (\modd k+1), \\ c \geqslant C}} \sum_{i \geqslant 0} \sum_{j \geqslant 0} (-1)^i \binom{l_C(c)-1+i}{i+j}\binom{i+j}{i}\binom{n}{a+i, c+j}.
    \end{equation}
    We can transform the right-hand side of Equation \ref{what_we_need_lemma}, completely analogously to the way we transformed the left-hand side into \ref{lhs_first_transformation}. After doing so, we get that the right-hand side also can be expressed as Expression \ref{lhs_first_transformation}, which finishes our proof.
\end{proof}

\end{document}
